\documentclass[11pt,a4paper]{article}

% ============================================================
% A11 - ANALISIS DE RIESGOS ROUTTECH
% Documento autocontenido y compatible con pdfLaTeX en Overleaf.
% ============================================================

\usepackage[utf8]{inputenc}
\usepackage[T1]{fontenc}
\usepackage[spanish,es-nodecimaldot]{babel}
\usepackage{tgheros}
\renewcommand{\familydefault}{\sfdefault}
\usepackage[a4paper,top=1.85cm,bottom=1.8cm,left=1.9cm,right=1.9cm,headheight=22pt]{geometry}
\usepackage{xcolor}
\usepackage{tabularx,array,colortbl,longtable}
\usepackage{fancyhdr}
\usepackage{lastpage}
\usepackage{microtype}
\usepackage{setspace}
\usepackage{titlesec}
\usepackage{hyperref}
\usepackage{xurl}
\usepackage{tikz}
\usepackage{ragged2e}
\usepackage{pdflscape}

% ---------- PALETA ROUTTECH ----------
\definecolor{RTForest}{HTML}{25463B}
\definecolor{RTTerracotta}{HTML}{B45F45}
\definecolor{RTGold}{HTML}{D5A64A}
\definecolor{RTCream}{HTML}{F7F2E8}
\definecolor{RTLight}{HTML}{EEF3F0}
\definecolor{RTGray}{HTML}{4D5652}
\definecolor{RTRed}{HTML}{A23B3B}
\definecolor{RTGreen}{HTML}{3F7D5A}
\definecolor{RTAmber}{HTML}{9C6B18}

% ---------- DATOS DEL DOCUMENTO ----------
\newcommand{\Proyecto}{RoutTech: Sistema de Movilidad Colaborativa para la Comunidad Universitaria de la UTEQ}
\newcommand{\Asignatura}{ISR-401 -- Ingeniería de Requisitos}
\newcommand{\Periodo}{2026--2027 PPA}
\newcommand{\CategoriaEtica}{Categoría B -- Datos personales}
\newcommand{\Repositorio}{https://github.com/AnderJM3/RoutTech_ISR401}
\newcommand{\FechaDocumento}{29 de julio de 2026}
\newcommand{\VersionDocumento}{2.0}
\newcommand{\Paralelo}{Cuarto nivel -- Software B}

% ---------- TABLAS ----------
\newcolumntype{Y}{>{\RaggedRight\arraybackslash}X}
\newcolumntype{C}[1]{>{\centering\arraybackslash}p{#1}}
\newcolumntype{L}[1]{>{\RaggedRight\arraybackslash}p{#1}}
\renewcommand{\arraystretch}{1.16}
\setlength{\tabcolsep}{4pt}

% ---------- ENCABEZADO Y PIE ----------
\pagestyle{fancy}
\fancyhf{}
\fancyhead[L]{\small\textbf{UTEQ | FCC | ISR-401}}
\fancyhead[R]{\small\textbf{A11 -- Análisis de riesgos}}
\fancyfoot[L]{\small Paquete ético RoutTech}
\fancyfoot[R]{\small Página \thepage\ de \pageref{LastPage}}
\renewcommand{\headrulewidth}{0.8pt}
\renewcommand{\footrulewidth}{0.4pt}
\renewcommand{\headrule}{\hbox to\headwidth{\color{RTTerracotta}\leaders\hrule height \headrulewidth\hfill}}
\renewcommand{\footrule}{\hbox to\headwidth{\color{RTGold}\leaders\hrule height \footrulewidth\hfill}}

% ---------- TITULOS ----------
\titleformat{\section}
  {\large\bfseries\color{RTForest}}
  {\thesection.}{0.45em}{}
  [\vspace{1pt}\color{RTGold}\titlerule]
\titlespacing*{\section}{0pt}{0.75em}{0.35em}

\hypersetup{
  colorlinks=true,
  linkcolor=RTForest,
  urlcolor=RTTerracotta,
  pdftitle={A11 Analisis de Riesgos RoutTech},
  pdfauthor={Equipo RoutTech}
}
\urlstyle{same}
\setlength{\parindent}{0pt}
\setlength{\parskip}{3pt}
\setstretch{1.08}

\begin{document}
\justifying

% ======================= PORTADA =======================
\thispagestyle{empty}
\begin{tikzpicture}[remember picture,overlay]
  \fill[RTForest] (current page.north west) rectangle ([yshift=-4.6cm]current page.north east);
  \fill[RTTerracotta] ([yshift=-4.6cm]current page.north west) rectangle ([yshift=-4.95cm]current page.north east);
\end{tikzpicture}

\vspace*{0.65cm}
\begin{center}
{\color{white}\Large\bfseries UNIVERSIDAD TÉCNICA ESTATAL DE QUEVEDO}

\vspace{1.75cm}
{\large\bfseries\color{RTForest} Facultad de Ciencias de la Computación\\[4pt]}
{\normalsize\color{RTGray} Carrera de Ingeniería de Software}

\vspace{0.8cm}

{\Huge\bfseries\color{RTForest} ANÁLISIS DE RIESGOS\\[8pt]}
{\Large\bfseries\color{RTTerracotta} A11 -- Medidas de mitigación}

\vspace{1.0cm}

\renewcommand{\arraystretch}{1.35}
\begin{tabularx}{0.90\textwidth}{>{\bfseries\color{RTForest}}L{4.0cm} Y}
\rowcolor{RTCream}
Proyecto & \Proyecto \\
Clasificación ética & \CategoriaEtica \\
Asignatura & \Asignatura \\
Periodo académico & \Periodo \\
Paralelo & \Paralelo \\
\end{tabularx}
\renewcommand{\arraystretch}{1.16}

\vfill
\begin{tabular}{rl}
\textbf{Documento:} & A11 -- Análisis de riesgos y medidas de mitigación \\
\textbf{Versión:} & \VersionDocumento \\
\textbf{Estado:} & Versión 2.0 para revisión \\
\textbf{Repositorio:} & \href{\Repositorio}{RoutTech\_ISR401} \\
\textbf{Fecha:} & \FechaDocumento
\end{tabular}

\vspace{0.75cm}
{\small Quevedo -- Los Ríos -- Ecuador}
\end{center}
\clearpage

% ======================= CONTENIDO =======================
\begin{landscape}
\thispagestyle{fancy}

\section{Metodología de valoración}

La matriz utiliza una evaluación cualitativa de \textbf{probabilidad}, \textbf{impacto} y \textbf{nivel de riesgo}. Cada riesgo incluye una medida concreta de prevención o mitigación.

\scriptsize
\begin{center}
\begin{tabularx}{0.96\linewidth}{L{2.5cm} Y Y Y}
\rowcolor{RTForest}
\textcolor{white}{\textbf{Nivel}} &
\textcolor{white}{\textbf{Probabilidad}} &
\textcolor{white}{\textbf{Impacto}} &
\textcolor{white}{\textbf{Tratamiento}}\\
\cellcolor{RTLight}\textcolor{RTGreen}{\textbf{Bajo}} & Poco probable. & Consecuencia menor y reversible. & Mantener controles y verificar.\\
\cellcolor{RTCream}\textcolor{RTAmber}{\textbf{Medio}} & Puede ocurrir. & Afectación moderada al estudio, datos o participantes. & Aplicar mitigación y supervisar.\\
\cellcolor{RTLight}\textcolor{RTRed}{\textbf{Alto}} & Probable o con exposición directa. & Daño relevante o incumplimiento ético. & Tratar antes de continuar.\\
\end{tabularx}
\end{center}

\section{Matriz de riesgos y medidas de mitigación}

\fontsize{7.35}{8.15}\selectfont
\renewcommand{\arraystretch}{1.00}
\setlength{\tabcolsep}{2.6pt}
\rowcolors{2}{RTLight}{white}
\begin{longtable}{C{0.75cm} L{2.6cm} L{4.9cm} C{1.35cm} C{1.35cm} C{1.35cm} L{8.0cm}}
\rowcolor{RTForest}
\textcolor{white}{\textbf{ID}} &
\textcolor{white}{\textbf{Categoría}} &
\textcolor{white}{\textbf{Riesgo}} &
\textcolor{white}{\textbf{Prob.}} &
\textcolor{white}{\textbf{Impacto}} &
\textcolor{white}{\textbf{Nivel}} &
\textcolor{white}{\textbf{Medida de prevención o mitigación}}\\
\endfirsthead
\rowcolor{RTForest}
\textcolor{white}{\textbf{ID}} &
\textcolor{white}{\textbf{Categoría}} &
\textcolor{white}{\textbf{Riesgo}} &
\textcolor{white}{\textbf{Prob.}} &
\textcolor{white}{\textbf{Impacto}} &
\textcolor{white}{\textbf{Nivel}} &
\textcolor{white}{\textbf{Medida de prevención o mitigación}}\\
\endhead
R-01 & Ético & Filtración de datos personales, consentimientos o transcripciones originales. & Media & Alto & Alto & Mantener los datos crudos fuera de GitHub, utilizar almacenamiento local cifrado y revisar los archivos antes de compartirlos.\\
R-02 & Ético & Publicación accidental de domicilios, recorridos o ubicaciones individuales. & Media & Alto & Alto & Generalizar origen y destino a sector, barrio, parroquia o campus; eliminar coordenadas y revisar cada evidencia antes de publicarla.\\
R-03 & Ético & Exposición de placas, rostros, voces o imágenes sin autorización. & Media & Alto & Alto & Solicitar autorizaciones separadas, aplicar desenfoque o recorte y excluir cualquier archivo identificable del repositorio.\\
R-04 & Ético & Uso de información para vigilancia académica, laboral o disciplinaria. & Baja & Alto & Medio & Prohibir usos ajenos al estudio, presentar resultados agregados y no elaborar perfiles individuales.\\
R-05 & Ético & Contacto no autorizado, hostigamiento o intercambio de datos entre participantes. & Baja & Alto & Medio & No compartir teléfonos ni correos; probar la comunicación del MVP únicamente con cuentas y datos sintéticos.\\
R-06 & Ético & Incomodidad, presión o malestar durante entrevistas y validaciones. & Media & Medio & Medio & Recordar el derecho a no responder, hacer pausas o retirarse sin consecuencias; detener la actividad ante cualquier malestar.\\
R-07 & Técnico & Pérdida, corrupción o eliminación accidental de evidencias. & Media & Medio & Medio & Mantener dos copias cifradas, verificar semanalmente su integridad y registrar los cambios de versión.\\
R-08 & Técnico & Carga accidental de archivos sensibles al repositorio público. & Media & Alto & Alto & Revisar cada commit, configurar \texttt{.gitignore} y retirar inmediatamente cualquier archivo publicado por error.\\
R-09 & Ético & Reidentificación de participantes mediante combinación de variables. & Media & Alto & Alto & Suprimir identificadores directos, reducir la precisión de variables y revisar el riesgo de reidentificación antes de divulgar datos.\\
R-10 & Reputacional & Afectación reputacional para participantes o para la UTEQ. & Baja & Alto & Medio & Evitar atribuir opiniones individuales, contextualizar los hallazgos y revisar el material institucional antes de difundirlo.\\
R-11 & Académico & Fabricación, alteración o interpretación no sustentada de datos. & Baja & Alto & Medio & Mantener trazabilidad entre evidencia, requisito y resultado; conservar los originales y documentar las transformaciones.\\
R-12 & Operativo & Suspensión o retiro de un aval institucional. & Baja & Alto & Medio & Detener la recolección afectada y obtener una nueva autorización antes de continuar.\\
R-13 & Operativo & Incumplimiento de plazos por reducción del equipo investigador. & Alta & Medio & Alto & Priorizar entregables, distribuir las tareas entre los dos integrantes y realizar seguimiento semanal del cronograma.\\
R-14 & Ético & Inclusión accidental de menores de edad. & Baja & Alto & Medio & Confirmar la mayoría de edad antes del consentimiento y excluir cualquier caso que no pueda verificarse.\\
R-15 & Técnico & Interpretación de una recomendación automática como decisión obligatoria. & Media & Medio & Medio & Mostrar los criterios, límites y advertencias de la recomendación, permitiendo rechazarla o modificarla.\\
\end{longtable}
\end{landscape}

\end{document}
